\documentclass[12pt]{article}
\usepackage[utf8]{inputenc}
\usepackage[T1]{fontenc}
\usepackage[margin=1in]{geometry}
\usepackage{setspace}
\usepackage{xcolor}
\usepackage[colorlinks=true, urlcolor=blue!60!black, linkcolor=blue!60!black]{hyperref}
\setlength{\parindent}{2em}
\setlength{\parskip}{1.25em}
\renewcommand{\baselinestretch}{1.0}
\title{\textit{Nixon v. Administrator of General Services}, 433 U.S. 425 (1977)}
\author{Case Brief}
\date{}
\begin{document}
\maketitle

\section{Case Metadata}

\begin{itemize}
  \item \textbf{Case Name:} \textit{Nixon v. Administrator of General Services}
  \item \textbf{Citation:} 433 U.S. 425 (1977) (\href{https://supreme.justia.com/cases/federal/us/433/425/}{Justia Primary Opinion}; \href{https://www.law.cornell.edu/constitution-conan/article-1/section-9/clause-3/bills-of-attainder}{Cornell LII Constitution Annotated})
  \item \textbf{Court:} Supreme Court of the United States (appeal from three-judge U.S. District Court for the District of Columbia, 408 F. Supp. 321 (1976))
  \item \textbf{Dates:} Argued April 20, 1977; Decided June 28, 1977
  \item \textbf{Vote:} 7--2 (affirming the District Court)
  \item \textbf{Majority Author:} Justice William J. Brennan, Jr. (joined in full by Justices Stewart, Marshall, and Stevens; joined in all but Part VII by Justice White; joined in all but Parts IV and V by Justice Powell; joined in Part VII by Justice Blackmun)
  \item \textbf{Concurring Opinions:} Justice Stevens (concurring); Justices White, Blackmun, and Powell (each concurring in part and concurring in the judgment)
  \item \textbf{Dissenting Opinions:} Chief Justice Burger; Justice Rehnquist
  \item \textbf{Chapter Placement:} Chapter 3: Federal Legislative Power (Bill of Attainder Clause, Art. I, \S~9, cl. 3; Separation of Powers; Executive Privilege)
\end{itemize}

\section{Facts and Procedural History}

Following the House Judiciary Committee's approval of articles of impeachment arising from the Watergate scandal, President Richard M. Nixon resigned on August 9, 1974. He directed federal archivists to ship his presidential materials---approximately 42 million pages of documents and 880 tape recordings---to California, but shipment was delayed after the Watergate Special Prosecutor advised President Gerald R. Ford of ongoing evidentiary needs. On September 8, 1974---the same day President Ford pardoned Nixon---General Services Administration (GSA) Administrator Arthur F. Sampson announced a depository agreement (the ``Nixon-Sampson agreement'') reciting that Nixon retained legal and equitable title to the materials. Under its terms, materials were stored under dual-key control near Nixon's California home; Nixon could withdraw original writings after three years and direct destruction of any tape recordings after September 1, 1979; and all tapes were to be destroyed upon his death or September 1, 1984, whichever came first (\href{https://supreme.justia.com/cases/federal/us/433/425/}{433 U.S. at 430--432}).

Ten days later, Congress moved to abrogate the agreement, enacting the Presidential Recordings and Materials Preservation Act (PRMPA), Title I of Pub. L. 93-526, 88 Stat. 1695, which President Ford signed on December 19, 1974. Title I directed the GSA Administrator to take complete possession and control of Nixon's presidential tapes and historical materials from January 20, 1969, through August 9, 1974; prohibited their destruction (\S~102(a)); made them available for judicial process subject to any legally or constitutionally based rights, defenses, or privileges (\S~102(b)); guaranteed Nixon ongoing access (\S~102(c)); and required regulations governing archival screening and eventual public access under seven statutory factors (\S~104(a)), including protecting legal privileges (\S~104(a)(5)) and returning purely personal and private materials to Nixon (\S~104(a)(7)). Section 105(c) authorized just compensation if the Act effected a taking of private property, while Title II (the Public Documents Act) created a national commission to study records standards for future Presidents and federal officials (\href{https://supreme.justia.com/cases/federal/us/433/425/}{433 U.S. at 433--436}).

On December 20, 1974, Nixon sued in the U.S. District Court for the District of Columbia seeking declaratory and injunctive relief on facial constitutional grounds. Because public-access regulations under \S~104 had not yet taken effect, a three-judge District Court limited review to the Act's custody and archival-screening provisions, rejected all constitutional claims, and dismissed the complaint (408 F. Supp. 321 (1976)). The Supreme Court noted probable jurisdiction (429 U.S. 976 (1976)) and affirmed.

\section{Issue(s) Presented}

\begin{enumerate}
  \item Does Title I of the PRMPA violate the Bill of Attainder Clause (U.S. Const. art. I, \S~9, cl. 3) by naming former President Nixon and overriding his depository agreement without a judicial trial?
  \item Does congressional regulation of a former President's records and archival screening by GSA violate Article II separation of powers or the qualified presidential privilege of confidentiality?
  \item Do GSA custody and archival screening facially violate Nixon's constitutional right of informational privacy or his First Amendment rights of political association and speech?
\end{enumerate}

\section{Holding}

\textbf{No on all counts.} The Court affirmed the dismissal and upheld the facial validity of Title I (\href{https://supreme.justia.com/cases/federal/us/433/425/}{433 U.S. at 426--428, 484}):

\begin{itemize}
  \item \textbf{Bill of Attainder:} Naming Nixon did not violate Art. I, \S~9, cl. 3 because he constituted a ``legitimate class of one'' and the Act did not inflict ``punishment'' under the Court's three-part \textbf{historical}, \textbf{functional}, and \textbf{motivational} test.
  \item \textbf{Separation of Powers:} Applying a pragmatic, functional disruption standard, the Act did not unduly prevent the Executive Branch from accomplishing its constitutionally assigned functions, especially because control remained within the Executive Branch (GSA) and both Presidents Ford and Carter supported the statute.
  \item \textbf{Presidential Privilege:} A former President has standing to assert the qualified privilege because it survives tenure for the benefit of the Republic, but the claim carries diminished weight when unsupported by the incumbent and yields here to minimal archival intrusion serving overriding public interests.
  \item \textbf{Privacy and First Amendment:} Nixon's legitimate privacy and associational interests in the small fraction of commingled personal and partisan records were outweighed by compelling governmental interests in preservation and safeguarded by discreet archival screening and statutory return provisions.
\end{itemize}

\section{Reasoning}

\subsection{Bill of Attainder Clause (Part VII)}

Justice Brennan rejected Nixon's attempt to read \textit{United States v. Brown}, 381 U.S. 437 (1965), as invalidating any law that singles out an individual for unwanted burdens. The Bill of Attainder Clause is not a variant of Equal Protection limiting Congress to ``legislating for the universe, or legislating only benefits, or not legislating at all'' (\href{https://supreme.justia.com/cases/federal/us/433/425/}{433 U.S. at 471}). Because all prior Presidents from Hoover through Johnson had already placed their papers in presidential libraries, and only Nixon had executed an agreement threatening imminent destruction of tapes upon his death or after five years, Nixon constituted a ``legitimate class of one,'' while Title II addressed future Presidents generally (\textit{id.} at 472).

Even assuming specificity was met, the Act did not ``inflict punishment'' under three distinct inquiries:
\begin{enumerate}
  \item \textbf{Historical Test:} Traditional attainder sanctions---death, imprisonment, banishment, vocational bars, and punitive confiscation of property---were absent. Section 105(c)'s provision for just compensation negated any claim of punitive confiscation (\textit{id.} at 473--475).
  \item \textbf{Functional Test:} Viewed by the type and severity of burdens imposed, the Act reasonably furthered legitimate, nonpunitive legislative goals: preserving evidence for Watergate criminal trials and safeguarding records of enduring historical value (\textit{id.} at 475--478).
  \item \textbf{Motivational Test:} Committee reports and floor debates reflected regulatory concern over the Nixon-Sampson agreement rather than legislative condemnation or intent to punish; indeed, the Act preserved Nixon's access (\S~102(c)), required return of personal items (\S~104(a)(7)), and granted expedited judicial review (\S~105(a)) (\textit{id.} at 478--484).
\end{enumerate}

\subsection{Separation of Powers and Presidential Privilege (Part IV)}

Rejecting an ``archaic view of the separation of powers as requiring three air-tight departments of government'' (\textit{id.} at 443), the Court held that the proper inquiry asks whether a statute prevents the Executive Branch from accomplishing its constitutionally assigned functions, and if so, whether an overriding need within Congress's authority justifies the impact. Here, GSA and its career archivists are Executive Branch employees, Congress received no privileged access beyond the general public, and both Presidents Ford and Carter endorsed the Act (\textit{id.} at 441--445).

On presidential privilege, the Court adopted the Solicitor General's view that confidentiality survives a President's tenure because advisers must rely on candor beyond a single term, but noted that the absence of incumbent support detracts from the weight of a former President's claim (\textit{id.} at 448--449). Analogizing archival screening to the \textit{in camera} judicial inspection approved in \textit{United States v. Nixon}, 418 U.S. 683 (1974), the Court concluded that screening by discreet professionals---coupled with \S~104(a)(5)'s privilege safeguards---imposed a minimal intrusion outweighed by Congress's substantial historical, legislative, and judicial-evidence interests (\textit{id.} at 450--455).

\subsection{Informational Privacy and First Amendment Rights (Parts V--VI)}

Citing \textit{Whalen v. Roe}, 429 U.S. 589 (1977), the Court recognized a legitimate expectation of privacy in purely personal materials (such as family correspondence and personal diaries), but balanced that interest against the public need. Because Nixon conceded seeing at most 200,000 of the 42 million pages and personal records were inextricably commingled with official business, archival screening---analogous to minimized electronic surveillance under Title III---was the least intrusive means to separate and return private items (\textit{id.} at 457--465). For similar reasons, any speculative chill on Nixon's First Amendment associational rights in partisan records was outweighed by compelling national interests (\textit{id.} at 465--468).

\section{Concurring and Dissenting Opinions}

\begin{itemize}
  \item \textbf{Justice Stevens (Concurring):} Joined the majority opinion but stressed two ``unmentioned'' facts decisive to his vote: Nixon's resignation under unique circumstances and his acceptance of a pardon. Without those facts, Stevens would have rejected single-person legislation aimed at a former President; he emphasized the case should not serve as precedent for future laws targeting an individual occupant rather than the Office (\href{https://supreme.justia.com/cases/federal/us/433/425/}{433 U.S. at 484--487}).
  \item \textbf{Justice White (Concurring in Part and in the Judgment):} Joined all but Part VII, relying on his \textit{Brown} dissent to find the Act nonpunitive, and insisted that purely private materials (such as diaries and Dictabelts) must be returned promptly regardless of historical interest (\textit{id.} at 487--491).
  \item \textbf{Justice Blackmun (Concurring in Part and in the Judgment):} Joined Part VII and largely agreed with Powell, but declined to treat the incumbent President's support as dispositive on separation of powers, warning against using the Act as a template during hostile presidential transitions (\textit{id.} at 491--492).
  \item \textbf{Justice Powell (Concurring in Part and in the Judgment):} Joined all but Parts IV and V; viewed Title I as unique emergency legislation and deemed President Ford's signature and President Carter's support dispositive of separation-of-powers and privilege claims, while cautioning that ``difficult constitutional questions lie ahead'' once public-access regulations issue (\textit{id.} at 492--504).
  \item \textbf{Chief Justice Burger (Dissenting):} Condemned the decision as bending settled constitutional principles under ``hydraulic pressure'' (\textit{id.} at 505). He argued Title I violated separation of powers by abandoning \textit{Humphrey's Executor} and \textit{Buckley v. Valeo} for a loose ``undue disruption'' test, invaded executive confidentiality and privacy, and functioned as an unconstitutional bill of attainder punishing a named former President (\textit{id.} at 504--545).
  \item \textbf{Justice Rehnquist (Dissenting):} Focused on separation of powers and executive confidentiality, warning that stripping former Presidents of control over their papers would chill candid advice and ``daily stand as a veritable sword of Damocles over every succeeding President and his advisers'' (\textit{id.} at 545, quoted at 494).
\end{itemize}

\section{Key Quotations}

\begin{quote}
``In short, while the Bill of Attainder Clause serves as an important `bulwark against tyranny,' \textit{United States v. Brown}, 381 U.S. at 443, it does not do so by limiting Congress to the choice of legislating for the universe, or legislating only benefits, or not legislating at all.''\\
--- Justice Brennan, \href{https://supreme.justia.com/cases/federal/us/433/425/}{433 U.S. at 471}
\end{quote}

\begin{quote}
``In short, appellant constituted a legitimate class of one, and this provides a basis for Congress' decision to proceed with dispatch with respect to his materials while accepting the status of his predecessors' papers and ordering the further consideration of generalized standards to govern his successors.''\\
--- Justice Brennan, \href{https://supreme.justia.com/cases/federal/us/433/425/}{433 U.S. at 472}
\end{quote}

\begin{quote}
``The Court, therefore, often has looked beyond mere historical experience and has applied a functional test of the existence of punishment, analyzing whether the law under challenge, viewed in terms of the type and severity of burdens imposed, reasonably can be said to further nonpunitive legislative purposes.''\\
--- Justice Brennan, \href{https://supreme.justia.com/cases/federal/us/433/425/}{433 U.S. at 475--476}
\end{quote}

\begin{quote}
``Rather, in determining whether the Act disrupts the proper balance between the coordinate branches, the proper inquiry focuses on the extent to which it prevents the Executive Branch from accomplishing its constitutionally assigned functions. \textit{United States v. Nixon}, 418 U.S. at 711--712. Only where the potential for disruption is present must we then determine whether that impact is justified by an overriding need to promote objectives within the constitutional authority of Congress.''\\
--- Justice Brennan, \href{https://supreme.justia.com/cases/federal/us/433/425/}{433 U.S. at 443}
\end{quote}

\section{Doctrinal Significance and Modern Legacy (Through 2026)}

\textit{Nixon v. Administrator of General Services} remains the foundational modern precedent in two distinct areas of constitutional law through 2026:

\begin{itemize}
  \item \textbf{Controlling Bill of Attainder Framework:} Reaffirmed in \href{https://supreme.justia.com/cases/federal/us/468/841/}{\textit{Selective Service System v. Minnesota Public Interest Research Group}}, 468 U.S. 841 (1984), and confirmed as the governing standard in \href{https://www.everycrsreport.com/files/2026-09-02_IF11293_c333cd4fdeb9bdf460eadc723c3b9ea6c1a30e27.html}{2026 Congressional Research Service guidance}, \textit{Nixon}'s three-part punishment test (historical, functional, motivational) has made Bill of Attainder challenges exceptionally difficult to win against regulatory or national-security statutes naming specific entities. Federal appellate courts invoked \textit{Nixon} to uphold appropriations bans in \href{https://caselaw.findlaw.com/court/us-2nd-circuit/1534979.html}{\textit{ACORN v. United States}}, 618 F.3d 125 (2d Cir. 2010), federal procurement bans targeting a named cybersecurity firm in \href{https://law.justia.com/cases/federal/appellate-courts/cadc/18-5176/18-5176-2018-11-30.html}{\textit{Kaspersky Lab, Inc. v. DHS}}, 909 F.3d 446 (D.C. Cir. 2018), and Congress's divest-or-ban mandate targeting ByteDance and TikTok in \href{https://caselaw.findlaw.com/court/us-dc-circuit/116750075.html}{\textit{TikTok Inc. v. Garland}}, 122 F.4th 930 (D.C. Cir. 2024) (holding forced divestment was a prophylactic national-security measure rather than punitive confiscation), \href{https://supreme.justia.com/cases/federal/us/604/24-656/}{\textit{aff'd on First Amendment grounds}}, 604 U.S. \_\_\_ (2025).
  \item \textbf{Presidential Records and Former-President Privilege:} Congress promptly enacted the Presidential Records Act of 1978, 44 U.S.C. \S\S~2201--2209, vesting public ownership of official presidential records in the United States. Decades later, \textit{Nixon}'s rule that an incumbent President's waiver of executive privilege outweighs a former President's objection proved decisive in \textit{Trump v. Thompson}, 20 F.4th 10 (D.C. Cir. 2021), \textit{stay denied}, 595 U.S. \_\_\_ (2022), permitting the National Archives to release January 6 White House records to Congress. At the same time, \textit{Nixon}'s functionalist ``undue disruption'' separation-of-powers test stands in growing tension with the modern Court's more formalist Article II jurisprudence, including \href{https://www.law.cornell.edu/supremecourt/text/23-939}{\textit{Trump v. United States}}, 603 U.S. 593 (2024).
\end{itemize}

\section{Pedagogical Critiques and Exam Hypotheticals}

\subsection{Critiques}

\begin{itemize}
  \item \textbf{Retreat from \textit{United States v. Brown}:} Critics note that Brennan's majority cited Justice White's \textit{Brown} dissent to narrow the Bill of Attainder Clause, transforming a structural bar on targeted legislation into a deferential inquiry where almost any plausible nonpunitive rationale rescues a named-target law.
  \item \textbf{Risk of Pretextual ``Classes of One'':} As Chief Justice Burger and Justice Stevens warned, allowing Congress to single out an unpopular individual whenever it pairs a targeted statute with a forward-looking study commission risks normalizing partisan targeting.
  \item \textbf{Functionalist Drift:} Relying on the incumbent President's consent to validate legislative control over a predecessor's records undervalues structural incentives in polarized transitions, a vulnerability Justice Blackmun foresaw.
\end{itemize}

\subsection{``Opposite Side'' Contrast: \textit{Foretich v. United States}}

To illustrate when a targeted statute \emph{fails} \textit{Nixon}'s three-part test, contrast \textit{Nixon} with \href{https://law.justia.com/cases/federal/appellate-courts/F3/351/1198/525998/}{\textit{Foretich v. United States}}, 351 F.3d 1198 (D.C. Cir. 2003). There, Congress passed the Elizabeth Morgan Act---narrowly drafted so it applied only to a single notorious D.C. custody dispute---stripping Dr. Eric Foretich of parental visitation rights without a judicial finding of abuse. Applying \textit{Nixon}, the D.C. Circuit invalidated the statute as an unconstitutional bill of attainder: unlike the PRMPA's compensated archival custody, the Act imposed a severe, permanent deprivation of parental rights disproportionate to any nonpunitive purpose (failing the functional test) and rested on legislative condemnation of Foretich as an abuser (failing the motivational test).

\subsection{Exam Hypotheticals}

\begin{enumerate}
  \item \textbf{Targeted Corporate Sanctions vs. Prophylactic Regulation:} Suppose Congress enacts a statute naming ``Apex AI Corp.'' and (a) prohibiting federal agencies from contracting with Apex due to documented foreign-intelligence vulnerabilities, while (b) forfeiting all of Apex's U.S. patents to the Treasury without compensation after floor speeches branding Apex's founders ``traitors.'' Under \textit{Nixon}, provision (a) likely survives all three punishment prongs as prophylactic procurement regulation (cf. \textit{Kaspersky}; \textit{TikTok}), whereas provision (b) fails both the historical test (uncompensated punitive confiscation) and the motivational/functional tests (cf. \textit{Foretich}).
  \item \textbf{Hostile-Transition Executive Privilege:} Suppose a newly elected President waives executive privilege over a predecessor's candid Oval Office national-security memos so a congressional committee can publish them, despite no active criminal trial or pending legislation. Students should analyze whether \textit{Nixon}'s diminished-weight rule for former Presidents (\href{https://supreme.justia.com/cases/federal/us/433/425/}{433 U.S. at 449}) controls, or whether Justice Blackmun's hostile-transition warning and modern Article II formalism (\textit{Trump v. Mazars}; \textit{Trump v. United States}) demand heightened judicial scrutiny.
\end{enumerate}

\end{document}
